\begin{tuctitlepage}
\begin{tuctitleorgunit}
Technische Universität Chemnitz \\
Fakultät für Textsatz \\
Professur \TeX-Prozessoren \\
\end{tuctitleorgunit}

\tuctitlelogo

\tuctitlethesistype{Studienarbeit}
\bigskip
\tuctitleauthor{Erika Musterfrau}
\bigskip
\begin{tuctitletableindent}{1}{0.5\linewidth}
Matrikel-Nr.: & 987654 \\
  geboren am: & 2.2.1996 \\
          in: & Chemnitz \\
 Studiengang: & Textsatz \\
\end{tuctitletableindent}

\vspace*{0pt plus 2fill}
\begin{tuctitletopic}
Die Bedeutung der Kommandosequenzen \\
\verb+\par+ und \verb+\everypar+
\end{tuctitletopic}
\vspace*{0pt plus 2fill}

\begin{tuctitletableindent}[\bfseries]{1}{0.5\linewidth}
Gutachter: & Prof.\,Dr.-Ing.\,habil. D. Ruck \\
Betreuer:  & Dipl.-Inf. Al Gorithmus \\
\end{tuctitletableindent}

\vspace*{0pt plus 1fill}

\tuctitleplacedate{Chemnitz, 1.4.2017}
\end{tuctitlepage}
